\section{Complete Mathematical Proofs}

All sums below are finite. Real-valued theorem statements include the rational
reference as a special case. None assumes statistical independence of keys,
values, or queries. The only independence used for sharpness is freedom to
choose each weight anywhere in its stated box.

\subsection{Convex Chord Bound and Continuous-Mass Sharpness}
\label{app:chord}

Suppose $\ell<u$ and $v\in[\ell,u]$. Put
$\alpha=(u-v)/(u-\ell)$ and $\beta=(v-\ell)/(u-\ell)$.
Then $\alpha,\beta\ge0$, $\alpha+\beta=1$, and
$v-t=\alpha(\ell-t)+\beta(u-t)$. The triangle inequality gives
\begin{equation}
 |v-t|\le\alpha|\ell-t|+\beta|u-t|.
\end{equation}
Summing and substituting $S=\sum_i v_i$ yields \cref{eq:chord}.
For $\ell=u$, all $v_i=\ell$, so the expression is exact. This proves
soundness without relying on a probabilistic interpretation of the moment.

For nonnegative real masses with total $n$ and first moment $S$, masses
$\lambda_\ell=(nu-S)/(u-\ell)$ and
$\lambda_u=(S-n\ell)/(u-\ell)$ at $\ell,u$ are feasible and attain
$C_{\rm ch}$. They are integers only for particular summaries. Choosing
weights $m+h\sgn(v-t)$ or $m-h\sgn(v-t)$ attains the corresponding
residual extreme for that endpoint distribution. This proves sharpness for
the continuous-mass relaxation, not generally for $n$ integer rows.

\subsection{Finite-Cardinality Extremality and Its Quantitative Gain}
\label{app:integer}

We prove \cref{thm:integer} constructively. The cases $\ell=u$, $S=n\ell$,
and $S=nu$ force a constant vector and agree with the stated formula.
Assume otherwise that $d=u-\ell>0$. Start with an arbitrary feasible
vector $v$. If two coordinates $x,y$ lie strictly inside $(\ell,u)$, define
\begin{equation}
 I=[\max(\ell-x,y-u),\ \min(u-x,y-\ell)].
\end{equation}
This is the feasible interval for replacing $(x,y)$ with
$(x+\delta,y-\delta)$. Its lower endpoint is strictly negative and its
upper endpoint strictly positive, so $0$ lies inside $I$. The function
$f(\delta)=|x+\delta-t|+|y-\delta-t|$ is convex. Writing $0$ as a convex
combination of the endpoints shows that at least one endpoint
$\delta_*$ has $f(\delta_*)\ge f(0)$. At that endpoint at least one of
the two coordinates equals $\ell$ or $u$. No previously fixed endpoint
coordinate changes, the sum is preserved, and the objective does not
decrease. Each operation reduces the number of interior coordinates by at
least one, so at most $n-1$ operations leave at most one interior coordinate.

If the resulting vector has one interior coordinate $x$, let $k$ be the
number of upper endpoints. Then
\begin{equation}
 S=ku+(n-k-1)\ell+x
   =n\ell+kd+(x-\ell).
\end{equation}
Since $0<x-\ell<d$, $k=\floor{(S-n\ell)/d}$ and
$x=x_*$. If no coordinate is interior, $r=(S-n\ell)/d$ is an integer.
For $r<n$ the formula with $x_*=\ell$ counts that endpoint once; for
$r=n$ the separate all-$u$ case applies. In every case the final objective
is exactly $C_{\rm int}(t)$. Because it is no smaller than the arbitrary
initial objective, it is an upper bound. The described final vector is
itself feasible, so the bound is attained. This proves the maximum equality
in \cref{eq:sharp}. Applying the chord inequality to the final vector
proves $C_{\rm int}\le C_{\rm ch}$.

For residual sharpness, \cref{lem:generic} supplies the bound. On a
maximizing value vector choose $w_i=m+h\sgn(v_i-t)$ to obtain
\begin{equation}
 \sum_iw_i(v_i-t)=m(S-nt)+h\sum_i|v_i-t|.
\end{equation}
These weights lie in $[a,b]$; at $v_i=t$ any allowed weight suffices.
Changing the sign before $h$ attains the lower endpoint. The denominator
may vanish if the box includes zero, but the residual extremality theorem
does not divide by it. Positivity is separately required for rounding.

To prove \cref{eq:improvement}, observe that the chord is exact at all
endpoint coordinates of the extremal vector. Its excess therefore comes
only from $x_*$. If $t\notin(\ell,u)$, absolute value is affine on
$[\ell,u]$ and the excess is zero. Otherwise let $g(x)$ be the chord
through $(\ell,|\ell-t|)$ and $(u,|u-t|)$. For $x\le t$,
\begin{equation}
 g(x)-|x-t|=\frac{2(x-\ell)(u-t)}{d},
\end{equation}
and for $x\ge t$ it equals $2(u-x)(t-\ell)/d$. The maximum occurs
at $x=t$ and is $2(t-\ell)(u-t)/d\le d/2$, because
$(t-\ell)(u-t)\le d^2/4$. Thus a single remainder coordinate accounts
for all the improvement and establishes its additive upper bound.

\subsection{Translation Invariance}

Replace $v_i,\ell,u,t,S$ by
$v_i+c,\ell+c,u+c,t+c,S+nc$. Then $u-\ell$, $S-n\ell$,
$nu-S$, and $S-nt$ are unchanged. The translated $r$ and $k$ are
unchanged, and $x_*$ becomes $x_*+c$. Every absolute difference in both
envelopes is unchanged. The exact residual is unchanged term by term.
Therefore both the chord and finite-cardinality residual enclosures are
translation invariant. This statement keeps weights fixed; it does not
claim invariance under score shifts in \cref{eq:profile}.

\subsection{Denominator Positivity, Cell Edges, and Proposal Soundness}

For each disjoint node $b$, a nonnegative weight box gives
$D_b\in[n_ba_b,n_bb_b]$. Summing the lower endpoints and obtaining a
strictly positive result establishes $D>0$. An exact opened leaf may
instead contribute a known positive denominator. A zero lower bound is
inconclusive even when the actual denominator is positive; a verifier may
refine, but cannot skip this premise.

For an arbitrary interval cell $J_y$ with independently open or closed
edges, the lower test is $\underline R(t_-)\ge0$ if the edge is
included and $\underline R(t_-)>0$ otherwise. The upper test similarly
uses $\overline R(t_+)\le0$ or $<0$. The identity
$R(t)=D(A-t)$ converts these to the exact membership conditions when
$D>0$. Every possible $A$ belongs to the same cell, so rounding is
decided independently of how the proposal was obtained.

BF16 has subnormal spacing $2^{-133}$, making the canonical zero cell
$[-2^{-134},2^{-134}]$. A normal power of two has unequal spacing to its
predecessor and successor, so replacing its true cell with a symmetric ULP
window can be unsound. At $16$, the predecessor is $255/16$ and the
successor is $129/8$; their midpoints with $16$ are $511/32$ and
$257/16$. The largest finite positive magnitude has odd significand and
an open overflow midpoint under this finite-only profile. Negative cells
are reflected from positive cells. The executable test enumerates all
finite canonical BF16 centers and checks these cell memberships directly
against exact rounding.

\subsection{Sound Summaries and Interval Queries}

Let a query coordinate satisfy $q_j\in[\underline q_j,\overline q_j]$
and a stored key coordinate satisfy
$k_{ij}\in[\underline k_j,\overline k_j]$ for every row. The product
belongs to the interval with endpoints the minimum and maximum of the four
endpoint products. Summing these intervals gives an enclosure $[L,U]$
of every dot product. Exact scaling and any sound bias enclosure may be
incorporated before applying the nondecreasing $\E$. Integer nearest-even
rounding is nondecreasing, as is $e^z$, hence $\E(L)\le\E(q^\top
k_i)\le\E(U)$. If a loose summary score interval exceeds the oracle
resource domain, the filter may open nodes or fall back to legal raw scores;
it must not infer that the reference invocation itself is invalid merely
from the loose box.

Summary merging follows directly from finite-sum associativity and extrema
over disjoint unions. For the positive/negative intersection, each positive
$v_i$ contributes in $[av_i,bv_i]$ and each negative $v_i$ in
$[bv_i,av_i]$. Summing yields $[aP+bM,bP+aM]$. The denominator interval,
signed interval scaling by $t$, and subtraction are sound operations.
Intersecting two sets containing the same residual preserves containment.
A disjoint result contradicts the premise that both sets contain it.

\section{Forest Invariants and Complexity Details}

Suppose the current completed roots cover consecutive disjoint ranges and
have strictly descending ranks. Appending a new full leaf creates a rank-zero
root at the right end. Merging the last two roots when they have the same
rank preserves their union, contiguity, and all summaries. Repeating the
carry produces exactly the set bits of the completed-leaf count. Induction
from the empty prefix proves root coverage. The partial tail, if nonempty,
starts at the end of the completed prefix and ends at $N$, completing the
partition.

If there are $m$ completed leaves and $p=\operatorname{popcount}(m)$
roots, exactly $m-p$ merges have occurred: each completion increases the
number of roots by one and each merge decreases it by one. In particular,
$m-p<m$ for $m>0$. One completion can still trigger
$\floor{\log_2m}+1$ carries. The summaries contain $O(d_k+d_v)$
scalars, and a full binary tree with $m$ leaves contains fewer than $2m$
nodes. With the partial tail, metadata storage is
$O((1+N/B)(d_k+d_v))$ scalar slots, in addition to
$O(N(d_k+d_v))$ raw data. These counts suppress integer/rational bit size.

For a query, replacing a node by disjoint children preserves both coverage
and sum of exact contributions. Each piece is marked unresolved or exact;
only unresolved pieces can be selected for refinement. Because descendants
of disjoint pieces are disjoint, fallback through unresolved subtrees never
revisits an exactified raw leaf. This is the operational condition underlying
\cref{thm:frontier}, not merely a property of the tree shape.

\section{Interval Interpreter and Transactional Composition}
\label{app:interpreter}

The interpreter claim is conditional on a specified computation graph and
sound primitive implementations. For any acyclic reference graph, prove
containment in topological order. Point inputs are enclosed exactly. An
interval affine map uses signed endpoint products and interval addition;
these contain the corresponding exact products and sums. Interval squaring
uses zero as its lower endpoint exactly when zero lies in the input interval.
For a nonnegative radicand interval, monotonicity of the square root allows
outward enclosures of its two endpoints. Reciprocal reverses endpoints on
a strictly positive interval. These facts preserve RMS normalization
containment after adding $\epsilon>0$.

The profile's exact point RMS calculation decides BF16 rounding by rational
square comparisons for $x/\sqrt r$; it need not approximate the irrational
number first. For non-point arguments, a square-root interval and outward
division suffice. The sigmoid factor $1/(1+\E(-x))$ is monotone, but
$x/(1+\E(-x))$ need not be. Multiplying enclosures of the two factors
therefore supplies the safe rule. Whenever a represented result is
materialized, monotonicity of the finite BF16 rounding map sends a real
interval to an interval of possible rounded values. If this interval is a
singleton, it equals the reference materialization. Overflow or an
unresolved numerical guard ends speculation.

Attention bounds contain all reference query/key/value combinations
represented by the current intervals. Their division requires a positive
denominator enclosure. Residual gates may replace a rounded attention
interval by an exact value using \cref{thm:gate}. Otherwise uncertainty may
propagate through further sound operations, until each live write is known
exactly or the attempt is rejected. This topological induction supplies the
interval-containment premise of \cref{thm:state}; it is not a formal
refinement proof of the Python interpreter.

Before execution, the caller retains the immutable original state.
Every proposed write is staged privately; a reference-equal scalar is
appended only to the staged representation. No reader may observe staging
before the final decision is certified. On rejection the staged object is
discarded, and dense evaluation starts from the original state, not from a
partially modified cache. On acceptance all persistent updates are
published as one logical transaction. Correct model/sequence attribution
and a complete state inventory are premises of this publication step.

For future-trace composition, let $\widehat F(C,x)$ be the filtered
transition. The numerical and write gates establish
$\widehat F(C,x)=F_R(C,x)$ on accepted invocations, and fallback establishes
the same equation on rejected invocations. For a supplied sequence
$x_0,\ldots,x_{T-1}$, the base states agree. If states at step $j$ agree,
the common input and the one-step equation give equal tokens and next
states. Induction proves all $T$ tokens and final states equal. For greedy
generation replace the system state by $(C,x)$ and use
$(C,x)\mapsto(y,(C',y))$; equal emitted tokens provide equal next inputs.
The statement concerns every finite legal continuation and contains no
accumulated numerical error or failure probability.

\section{Outward Exponential Enclosures and E24 Rounding}
\label{app:lattice}

Let $\Delta=2^{-24}$ and define
$Q_\Delta(x)=\Delta\RNE_{\mathbb Z}(x/\Delta)$ on nonnegative real
inputs. Suppose a sound exponential computation gives
$0\le L\le e^z\le U$. The quantized weight satisfies
\begin{equation}
 \E(z)\in[Q_\Delta(L),Q_\Delta(U)]
 \subseteq[\max(0,L-\Delta/2),U+\Delta/2].
 \label{eq:lattice}
\end{equation}
Thus rounding the outward enclosure endpoints to the target lattice is
at least as strong as adding a global half-quantum error, provided the
endpoint quantization itself is exact.

\begin{proof}
Nearest-even rounding is nondecreasing. To see this directly, suppose
$x\le y$ but their rounded lattice values satisfy $a>b$. Nearest-point
optimality implies $|x-a|\le|x-b|$ and $|y-b|\le|y-a|$, giving
$x\ge(a+b)/2\ge y$. Hence $x=y$, whose deterministic rounding cannot
give different outputs, a contradiction. Applying this monotonicity to
$L\le e^z\le U$ proves the first inclusion. Every nonnegative input is
within $\Delta/2$ of a nonnegative lattice point, so
$|Q_\Delta(x)-x|\le\Delta/2$; the second inclusion follows.
\end{proof}

The implementation uses exact power-of-two scaling, integer parity, and
nearest-even tie handling for admitted finite binary64 endpoints, rejecting
overflow or invalid inputs. The generic monotone endpoint implication is
covered by the compiled \code{StateCut.monotone\_weight\_bounds}; the
binary64 implementation and exponential enclosure are separately tested,
and are not formally refined to that theorem. An exact negative-tail case
uses $z\le-25$: since $e>2$, $e^z\le e^{-25}<2^{-25}=\Delta/2$,
so $\E(z)=0$. Neither construction applies a common score shift.

\section{A Backend Bridge Is an Additional Obligation}
\label{app:bridge}

\begin{proposition}[Inward-shifted bridge]
Let $D>0$, $A=N/D$, $\varepsilon\ge0$, and suppose a backend's
pre-materialization scalar $B$ satisfies $|B-A|\le\varepsilon$.
If residual bounds establish
\begin{equation}
 R(t_-+\varepsilon)\ge0,\qquad
 R(t_+-\varepsilon)\le0,
\end{equation}
with strict signs for excluded cell edges, then $B\in J_y$ and its
materialization in the same rounding format is $y$.
\end{proposition}
\begin{proof}
The residual signs give $A\ge t_-+\varepsilon$ and
$A\le t_+-\varepsilon$. Since $B\ge A-\varepsilon$ and
$B\le A+\varepsilon$, $t_-\le B\le t_+$. A strict first or second
inequality gives the corresponding strict endpoint condition. These are
exactly the requirements for membership in $J_y$.
\end{proof}

The result does not construct $\varepsilon$. In particular it cannot be
instantiated by an empirical tolerance, a few matching output tensors, or
a theorem about a different exponential. If the inward-shifted cell is
empty the sufficient conditions cannot pass, which correctly reflects the
insufficiency of that error bound for exact rounding. A state-level bridge
must cover every numerical operation contributing to a live write, not
only the final attention quotient.

\section{Reproducibility and Artifact Map}
\label{app:reproduce}

The repository retains raw machine-readable measurements and their logs.
The manuscript's generated tables contain no manually entered measured
numbers; their input paths and SHA-256 hashes appear in
\code{paper/generated/input\_manifest.json}. Current GH200 measurements
are retained under \code{results/gh200/}; historical release snapshots do
not substitute for a successful current execution.

\begin{center}
\begin{tabular}{p{0.47\textwidth}p{0.46\textwidth}}
\toprule
Artifact & Responsibility \\
\midrule
\code{src/statecut/arithmetic.py} & Exact rounding, elementary enclosures, greedy tie rule \\
\code{src/statecut/residual.py} & Chord and finite-count residual predicates \\
\code{src/statecut/forest.py} & Persistent summary construction and refinement tree \\
\code{src/statecut/tree\_attention.py} & Attention certificates and exact raw-leaf fallback \\
\code{src/statecut/tree\_model.py} & Exact cuts and staged persistent-write checks \\
\code{scripts/run\_v2\_experiments.py} & Paired logical-read and complete-state experiments \\
\code{scripts/capture\_pretrained.py},\newline \code{scripts/audit\_capture.py} & Backend-preserving capture and offline audit \\
\code{lean/StateCut/} & Theorem sources with a pinned Lean/Mathlib environment \\
\code{docs/LITERATURE\_AUDIT.md} & Primary-source comparison and novelty scope \\
\code{paper/} & ICML preprint source, official style, bibliography, and PDF \\
\bottomrule
\end{tabular}
\end{center}

From the repository root, install the test dependencies and run:
\begin{verbatim}
python -m pip install -e '.[test]'
python -m pytest -q
python scripts/run_v2_experiments.py \
  --output results/gh200/algorithm/v2_experiments.json
bash scripts/check_lean.sh
make -C paper
\end{verbatim}
The Lean check must return success after compiling the actual declarations;
source scans for forbidden placeholders are supporting checks only. CUDA
build and capture commands, compiler flags, model revision, dtype, attention
implementation, prompt token IDs, and device metadata belong in the
corresponding run record. Capture tests can skip if optional Torch or NumPy
packages are missing; both pass and skip counts must be inspected.

The official ICML 2026 style is vendored without modification from the
conference's author-instruction archive, whose source and hash are in
\code{paper/template\_source.json}. The manuscript uses the style's
\code{preprint} option.

\section{Pretrained CUDA Domain Accounting}
\label{app:domains}
% Generated by render_results.py; do not hand-edit.
\begin{center}
\begin{tabular}{rrrr}
\toprule
Block size & Unsupported weights & Positive mass unproved & Valid positive mass \\
\midrule
1 & 346 & 4 & 10,558 \\
16 & 415 & 242 & 10,251 \\
64 & 462 & 234 & 10,212 \\
128 & 524 & 224 & 10,160 \\
\bottomrule
\end{tabular}
\end{center}
The columns partition the query heads within each configuration. Unsupported weight enclosures are rejected before the output predicate; a nonpositive lower denominator bound does not establish that the actual denominator is nonpositive. The two proposal methods share these counts. Accepted-head counts use the full population, including rejected domains, rather than silently excluding unsupported cases. CUDA admits a narrower positive-score domain than the rational oracle. The capture audit requires the exact query scale $1/8$, as used by these head-dimension-64 models; arbitrary rounded query scaling is outside this audit's contract.


\section{Compiled Formal Inventory}
\label{app:formal}
% Generated by render_results.py; do not hand-edit.
The build uses \code{leanprover/lean4:v4.19.0}, with Mathlib at \nolinkurl{c44e0c8ee63ca166450922a373c7409c5d26b00b}. The audit counts source theorem declarations separately from generated logical helpers. It records 14 compiler runtime extraction auxiliaries outside the proof inventory; these are not used to replace kernel checking. The reported counts and names come from the build audit JSON.

\begin{center}

\begin{tabular}{p{0.29\textwidth}p{0.65\textwidth}}

\toprule

Claim family & Compiled declarations in namespace \code{StateCut} \\

\midrule

Moment and residual bounds & \code{abs\_chord\_sum}, \code{moment\_residual\_sound}, \code{finite\_abs\_sum\_bound}, \code{finite\_residual\_bound} \\

Finite-count strength and feasibility & \code{finite\_envelope\_le\_chord}, \code{finite\_envelope\_attained}, \code{floor\_remainder\_contract} \\

Rounding-cell composition & \code{residual\_to\_cell}, \code{moment\_to\_cell}, \code{finite\_to\_cell} \\

Conditional numerical bridge & \code{backend\_bridge\_gate}, \code{backend\_bridge\_cell} \\

Disjoint frontier sums & \code{disjoint\_split\_sum}, \code{frontier\_sum\_bounds} \\

Persistent writes and traces & \code{singleton\_writes}, \code{checked\_write\_frontier}, \code{write\_frontier\_future} \\

\bottomrule

\end{tabular}

\end{center}

The finite-count theorem takes an integer endpoint count and a remainder satisfying the summary equation. The separate mathematical floor/remainder theorem supplies those conditions in the nonconstant, non-upper-endpoint case. Constant cases are proved separately. Neither result establishes that Python's runtime executes the specified mathematics; the exact-rational tests check that implementation separately. The cell theorems take a correct rounding-cell contract as a premise and do not prove the full BF16 encoding/cell generator. The frontier results concern disjoint sums; they do not formally verify the Python persistent forest or its memory-access counters.


\section{Interpretation of Negative Results}

An exact certificate is useful only if its premises can be established
cheaply and it accepts often enough for the intended workload. Failure of
an interval gate does not show that the proposal is wrong; it shows that
the current enclosure cannot prove it right. Full raw fallback with exact
state parity is therefore a valid correctness outcome and a negative
performance outcome. Similarly, a zero-error device test is finite
conformance evidence, while a counterexample can falsify a universal
implementation claim immediately. The present artifact preserves both
types of result.
